\documentclass[11pt]{article}
\usepackage[margin=2.6cm]{geometry}
\usepackage{amsmath,amssymb}
\usepackage{booktabs}
\usepackage{hyperref}
\usepackage{parskip}
\usepackage{xcolor}
\usepackage{listings}
\usepackage{enumitem}

\lstset{
  basicstyle=\ttfamily\small,
  columns=fullflexible,
  breaklines=true,
  frame=single,
  language=R
}

\newcommand{\code}[1]{\texttt{#1}}
\newcommand{\file}[1]{\textcolor{black!60}{\texttt{#1}}}

\title{Derivative (Curvature) Regularization for ODE Model Fits\\[2pt]
\large A step-by-step guide to the code and the math}
\author{}
\date{Reference: \file{scripts/derivative\_constraint\_relay.R}\\(machinery shared verbatim with \file{scripts/derivative\_constraint.R})}

\begin{document}
\maketitle

This note explains, from first principles, what \code{curvatureConstraint()}
does and why it is built the way it is. It assumes you know a little R, but
nothing about ODE sensitivity analysis, finite differences, or optimization
theory -- every piece of math is derived, not just stated.

\tableofcontents

\section{The problem, in plain terms}

A dynamic model describes how a set of quantities (``states'') change over
time via a differential equation $\dot{X} = f(X,\theta)$, where $\theta$ is
a vector of unknown parameters (rate constants, initial fractions, ...).
Fitting the model means choosing $\theta$ so that the model's predictions
$Y(t) = g(X(t))$ -- the \emph{observable} quantities, some function of the
states -- match measured data as closely as possible.

The trouble: when the data is sparse, noisy, or only covers \emph{some} of
the states, many different $\theta$ can fit the data almost equally well.
This is called \textbf{practical non-identifiability} (as opposed to
\emph{structural} non-identifiability, which is a property of the model
equations alone, true even with infinite perfect data -- practical
non-identifiability is a property of what you actually measured). Among all
the equally-good-fitting $\theta$'s, nothing in the data rules out ones that
produce biologically implausible trajectories: sharp jumps, ringing, wild
excursions in the parts of the model that were never directly observed.
A plain maximum-likelihood fit has no way to prefer the plausible member of
that family over an implausible one -- it only sees the likelihood, and by
construction they all score about the same.

\textbf{The fix explored here:} add an extra term to the fitting objective
that specifically penalizes \emph{sharp corners in time} (formally: high
second time-derivative, i.e.\ curvature), so the optimizer is nudged toward
smoother, more plausible reconstructions -- without needing to know in
advance which specific parameter is the culprit.

\section{Minimal background}

\subsection{States, observables, parameters}
A state trajectory $X(t;\theta)$ solves $\dot X = f(X,\theta)$ from some
initial condition. An observable is $Y(t;\theta) = g(X(t;\theta))$ -- in the
relay model, e.g.\ $\code{obsTF} = \code{scale\_TF}\cdot\code{TFPnuc}$, a
scalar multiple of one state. Parameters are almost always fit on a
$\log$-scale (i.e.\ $\theta = \log(\text{rate constant})$) so that the
optimizer's unconstrained search automatically respects that a rate
constant must stay positive.

\subsection{The objective, and why it needs a gradient and a Hessian}
Fitting means minimizing an objective $\mathrm{obj}(\theta)$ -- essentially
a (noise-weighted) sum of squared differences between model and data. Almost
every practical optimizer needs, in addition to the value
$\mathrm{obj}(\theta)$ itself:
\begin{itemize}[nosep]
  \item the \textbf{gradient} $\nabla_\theta\,\mathrm{obj}$ -- which
  direction in parameter space improves the fit fastest;
  \item the \textbf{Hessian} $\nabla^2_\theta\,\mathrm{obj}$ -- how much the
  objective curves in parameter space, used to size a good step.
\end{itemize}

\textbf{A warning about the word ``curvature'':} this document uses
\emph{curvature} in two completely different senses, and they must not be
confused:
\begin{enumerate}[nosep]
  \item \textbf{Curvature in time} -- how sharply a trajectory $Y(t)$ bends
  as $t$ increases. This is what the new penalty targets, and it is a
  property of the \emph{predicted trajectory}, nothing to do with $\theta$
  directly.
  \item \textbf{Curvature in parameter space} -- the Hessian
  $\nabla^2_\theta\,\mathrm{obj}$, how the \emph{objective} bends as you
  wiggle $\theta$. This is what the optimizer needs to take good steps.
\end{enumerate}
Confusingly, both matter here: we build a penalty on time-curvature (\S3-8),
and then we still need that penalty's own gradient/Hessian \emph{in
parameter space} (\S9-10) so the optimizer can act on it.

\subsection{Sensitivities: $\partial Y/\partial\theta$, computed exactly}
To get an exact gradient of anything built from $Y(t;\theta)$, you need
$\partial Y(t;\theta)/\partial\theta$ -- how much the prediction at time $t$
would change if you nudged parameter $\theta_k$. dMod does not approximate
this by finite-differencing $\theta$ (nudge $\theta_k$, re-solve the ODE,
see what changed); instead it integrates extra ``sensitivity'' equations
\emph{alongside} the model equations, giving the exact derivative at every
time point as a direct by-product of the same ODE solve. You can treat this
as a black box for the rest of this document: whenever a prediction is
computed, an array
\[
  S(t) \;=\; \frac{\partial Y(t;\theta)}{\partial\theta} \quad
  \text{(shape [time, quantity, }\theta\text{])}
\]
comes along with it, attached as \code{attr(prediction, "deriv")}, exactly
and at no extra approximation cost.

\section{What ``curvature'' means here}

Think of $Y(t)$ as a position. Its first derivative $\dot Y(t)$ is a
velocity -- how fast it's changing. Its second derivative $\ddot Y(t)$ is an
acceleration -- how fast the \emph{velocity} is changing, i.e.\ how sharply
the curve is bending.

\begin{itemize}[nosep]
  \item A straight line has zero curvature everywhere, even if it is moving
  fast (large $\dot Y$, but constant slope means $\ddot Y = 0$).
  \item A smooth, gently-curving rise-and-fall (like a bell curve) has
  small, smoothly-varying curvature.
  \item A sharp corner -- a kink where the slope changes abruptly -- has, in
  the limit, curvature blowing up to infinity right at the kink.
\end{itemize}

This is exactly why the \emph{second} derivative, not the first, is the
right quantity to penalize (as discussed when this was first proposed):
penalizing $\dot Y$ directly would also punish anything that moves fast,
including genuinely fast-but-smooth dynamics or oscillation. Penalizing
$\ddot Y$ specifically targets \emph{corners}, and leaves fast-but-smooth
motion alone.

\section{Approximating $\ddot Y$: finite differences from scratch}

We only ever have $Y$ at a discrete set of times, so $\ddot Y(t)$ must be
\emph{approximated} from nearby values. The derivation, from Taylor's
theorem:
\begin{align}
  Y(t+h) &= Y(t) + h\,\dot Y(t) + \tfrac{h^2}{2}\,\ddot Y(t) + \tfrac{h^3}{6}\,\dddot Y(t) + O(h^4) \\
  Y(t-h) &= Y(t) - h\,\dot Y(t) + \tfrac{h^2}{2}\,\ddot Y(t) - \tfrac{h^3}{6}\,\dddot Y(t) + O(h^4)
\end{align}
Adding these two lines, the odd-order terms ($\dot Y$, $\dddot Y$) cancel
exactly:
\begin{equation}
  Y(t+h) + Y(t-h) \;=\; 2\,Y(t) + h^2\,\ddot Y(t) + O(h^4)
\end{equation}
Rearranging for $\ddot Y(t)$:
\begin{equation}
  \boxed{\ \ddot Y(t) \;\approx\; \frac{Y(t+h) - 2\,Y(t) + Y(t-h)}{h^2}\ },
  \qquad \text{error } O(h^2).
  \label{eq:central2diff}
\end{equation}
This is the \emph{central second difference} -- and it is exactly
\code{secondDiffTime()} (\file{derivative\_constraint\_relay.R:255-263}):
\begin{lstlisting}
secondDiffTime <- function(arr, h) {
  n <- dim(arr)[1]
  lo <- 1:(n - 2); mid <- 2:(n - 1); hi <- 3:n
  if (length(dim(arr)) == 2) {
    (arr[hi, , drop = FALSE] - 2 * arr[mid, , drop = FALSE] + arr[lo, , drop = FALSE]) / h^2
  } else {
    (arr[hi, , , drop = FALSE] - 2 * arr[mid, , , drop = FALSE] + arr[lo, , , drop = FALSE]) / h^2
  }
}
\end{lstlisting}
\code{lo}, \code{mid}, \code{hi} are exactly $t-h$, $t$, $t+h$ in
Equation~\eqref{eq:central2diff}, shifted by one index each. Because the
formula needs a point on \emph{both sides}, it can only be evaluated at
\emph{interior} grid points -- the first and last time point have no
`before'/`after' neighbor, so the result array always has 2 fewer rows than
the input (\code{n-2} instead of \code{n}). The function also handles a
3-dimensional array (\code{[time, quantity, }$\theta${]}) the same way,
differencing along the time axis only -- this is what makes it reusable for
the sensitivity array too (\S9).

This formula assumes the grid spacing $h$ is the \emph{same} everywhere.
That is why \code{curvatureConstraint()} checks it explicitly:
\begin{lstlisting}
stopifnot("curvatureConstraint() assumes a uniform time grid" =
            max(abs(hs - hs[1])) < 1e-8 * hs[1])
\end{lstlisting}
and why the fitting grid is deliberately built as an evenly-spaced sequence,
\code{seq(0, 80, by = 0.1)} (\S13).

\section{From ``the trajectory is rough'' to ``a number to minimize''}

Equation~\eqref{eq:central2diff} gives one curvature value per quantity $j$
(a state or observable) per interior time point $t_i$:
\[
  \mathrm{curv}_j(t_i;\theta) \;=\; \frac{Y_j(t_{i+1};\theta) - 2\,Y_j(t_i;\theta) + Y_j(t_{i-1};\theta)}{h^2}.
\]
To turn ``how rough is quantity $j$'s whole trajectory'' into a single
number, average its \emph{squared} curvature over all interior time points
(squared, not absolute value, for the same reason least-squares fitting uses
squared residuals rather than absolute residuals: $x^2$ is smooth
everywhere, including at $x=0$, so the whole penalty stays differentiable):
\[
  \overline{\mathrm{curv}_j^2}(\theta) \;=\; \frac{1}{n_t}\sum_{i} \mathrm{curv}_j(t_i;\theta)^2,
  \qquad n_t = \text{number of interior time points}.
\]

\section{Different scales: self-normalizing weights}

Different quantities live on very different numeric scales -- a fast
receptor pulse decaying from $0.8$ to near $0$, versus a slowly accumulating
nuclear pool -- and their \emph{natural} curvature differs accordingly.
Summing raw $\overline{\mathrm{curv}_j^2}$ across quantities would let
whichever one happens to have the biggest numbers dominate the penalty, for
reasons that have nothing to do with whether it is being fit
\emph{implausibly}.

The fix: give every quantity its own reference scale, measured once, at a
reference fit called \code{best} (an ordinary, unregularized fit):
\[
  \mathrm{ref}_j \;=\; \overline{\mathrm{curv}_j^2}(\code{best}),
  \qquad
  \mathrm{weight}_j \;=\; \frac{1}{\mathrm{ref}_j + \varepsilon}.
\]
\begin{lstlisting}
ref <- colMeans(curv_best^2)
eps_curv <- eps_curv_frac * max(ref)
weight <- 1 / (ref + eps_curv)
\end{lstlisting}
Why this works: evaluated \emph{at the reference point itself},
$\mathrm{weight}_j \cdot \overline{\mathrm{curv}_j^2}(\code{best}) =
\mathrm{ref}_j/(\mathrm{ref}_j+\varepsilon) \approx 1$ for every quantity
$j$ (as long as $\mathrm{ref}_j \gg \varepsilon$) -- regardless of that
quantity's raw numeric scale. So every state and observable starts out
contributing \emph{about the same amount} to the total penalty; only
\emph{deviations} from the reference curvature (getting rougher, in either
direction) actually move the penalty.

$\varepsilon = \code{eps\_curv\_frac} \times \max_j(\mathrm{ref}_j)$ is a
tiny floor (a \emph{proportional} one, not a fixed absolute number, so it
scales with whatever units the model happens to use) that prevents a
division blow-up for a quantity that happens to be almost perfectly flat at
\code{best} (where $\mathrm{ref}_j \approx 0$).

\section{One interpretable knob: \code{rel\_strength}}

We still need an overall scale for how hard to enforce smoothness. Rather
than an arbitrary absolute number, it is defined as a \emph{fraction of the
model's own likelihood value at the reference fit}:
\[
  \mathrm{strength\_curv} \;=\; \mathrm{rel\_strength} \times \bigl|\mathrm{obj}(\code{best})\bigr|.
\]
\begin{lstlisting}
strength_curv <- rel_strength * abs(out_best$value)
\end{lstlisting}

\textbf{Why \code{abs()} -- a real bug, worth understanding.} The fitting
objective's value need not be positive: it is a full Gaussian negative
log-likelihood, and for small measurement noise $\sigma$ its normalization
term ($\propto \log\sigma^2$) can dominate and make the \emph{total} value
negative (it does, here: $\mathrm{obj}(\code{best})\approx-54.7$). Without
\code{abs()}, \code{strength\_curv} would inherit that sign and turn
\emph{negative} -- and a negative-strength penalty \emph{rewards} roughness
instead of discouraging it. This was caught empirically while building this
script (Section~14 revisits the general lesson it points to), which is
exactly the kind of subtlety that is easy to miss on paper and cheap to
catch by testing.

\textbf{Why this choice of normalization is interpretable.} Averaging
(rather than summing) over both quantities and time points --
$C = \mathrm{strength\_curv}/(n_{\mathrm{cols}}\cdot n_t)$ -- means that,
evaluated at \code{best} itself,
\[
  R(\code{best}) \;=\; C\sum_j \mathrm{weight}_j \sum_i \mathrm{curv}_j(t_i;\code{best})^2
  \;\approx\; \frac{\mathrm{strength\_curv}}{n_{\mathrm{cols}}}\sum_j 1
  \;=\; \mathrm{strength\_curv},
\]
using the $\approx 1$ result from \S6. So \code{rel\_strength = 0.1} means,
concretely: \emph{at the reference fit, the smoothness penalty is worth
about 10\% of the data-fit term} -- a directly interpretable statement,
rather than an arbitrary number with no intrinsic scale.

\section{The penalty function, assembled}

Putting \S5-7 together, the full penalty is
\begin{equation}
  R(\theta) \;=\; C \sum_{j=1}^{n_{\mathrm{cols}}} \mathrm{weight}_j
  \sum_{i=1}^{n_t} \mathrm{curv}_j(t_i;\theta)^2,
  \qquad C = \frac{\mathrm{strength\_curv}}{n_{\mathrm{cols}}\, n_t}.
  \label{eq:penalty}
\end{equation}
\begin{lstlisting}
pen_value <- C * sum(vapply(seq_len(n_cols), function(j)
  weight[j] * sum(curv[, j]^2), numeric(1)))
\end{lstlisting}
Here \code{curv} is \code{secondDiffTime(Y, hs)} applied to the trajectory
\emph{at the current $\theta$} being tried by the optimizer -- not at
\code{best}. \code{best} only ever supplies the fixed reference weights.

\section{The gradient: exact, not finite-differenced}

The optimizer needs $\partial R/\partial\theta_k$. By the ordinary chain
rule applied to Equation~\eqref{eq:penalty}:
\begin{equation}
  \frac{\partial R}{\partial\theta_k} \;=\; 2C\sum_j \mathrm{weight}_j \sum_i
  \mathrm{curv}_j(t_i;\theta)\;\frac{\partial\, \mathrm{curv}_j(t_i;\theta)}{\partial\theta_k}.
  \label{eq:chainrule}
\end{equation}
Everything hinges on that last factor. Write out
$\mathrm{curv}_j(t_i;\theta) = [Y_j(t_{i+1};\theta) - 2Y_j(t_i;\theta) +
Y_j(t_{i-1};\theta)]/h^2$: it is a \emph{fixed linear combination} (with
coefficients $1,-2,1$, divided by $h^2$ -- none of which depend on $\theta$)
of $Y_j$ evaluated at three times. Because differentiation is linear
--$\ \partial(aU+bV)/\partial\theta_k = a\,\partial U/\partial\theta_k +
b\,\partial V/\partial\theta_k$ for constants $a,b$ -- differentiating this
fixed combination with respect to $\theta_k$ just applies the \emph{same}
combination to the derivatives:
\begin{equation}
  \frac{\partial\,\mathrm{curv}_j(t_i;\theta)}{\partial\theta_k}
  \;=\; \frac{S_j(t_{i+1};\theta)_k - 2\,S_j(t_i;\theta)_k + S_j(t_{i-1};\theta)_k}{h^2},
  \qquad S_j(t;\theta)_k \equiv \frac{\partial Y_j(t;\theta)}{\partial\theta_k}.
  \label{eq:commute}
\end{equation}
This is \emph{the same finite-difference operator}, applied not to $Y$ but
to the sensitivity $S$ -- because the time grid $h$ never depends on
$\theta$, differencing in time and differentiating in $\theta$
\textbf{commute exactly}. There is no approximation introduced by this
step beyond the one already made in \S4 for the trajectory itself. This is
the central trick of the whole construction, and it is why the code can
reuse \code{secondDiffTime()} completely unchanged:
\begin{lstlisting}
dY    <- attr(pred, "deriv")[, cols, , drop = FALSE]   # S(t), exact, from dMod
dcurv <- secondDiffTime(dY, hs)                        # exactly Eq. (7)
\end{lstlisting}
Substituting Equation~\eqref{eq:commute} into Equation~\eqref{eq:chainrule}
and writing it as a matrix product (summing over $i$ for every $k$ at
once, instead of looping):
\begin{lstlisting}
for (j in seq_len(n_cols)) {
  Jc <- dcurv[, j, ]                 # matrix: rows = time, cols = theta
  dim(Jc) <- c(dim(dcurv)[1], length(theta))
  pen_gradient <- pen_gradient + 2 * C * weight[j] * as.vector(curv[, j] %*% Jc)
  ...
}
\end{lstlisting}
\code{curv[, j] \%*\% Jc} computes, in one matrix-vector product,
$\sum_i \mathrm{curv}_j(t_i)\cdot\mathrm{Jc}_{i,k}$ for every $\theta_k$
simultaneously -- exactly the sum in Equation~\eqref{eq:chainrule}. The
$\theta$-derivative of the whole penalty is therefore \textbf{exact given
the time discretization} -- $\theta$ itself is never finite-differenced,
only time is (and only once, in \S4).

\section{The Hessian: a Gauss--Newton approximation}

The optimizer also wants $\partial^2 R/\partial\theta_k\partial\theta_l$.
Differentiating Equation~\eqref{eq:chainrule} once more with respect to
$\theta_l$ (product rule on the two $\theta$-dependent factors) gives
\[
  \frac{\partial^2 R}{\partial\theta_k\partial\theta_l} =
  2C\sum_j w_j \sum_i \left[
    \frac{\partial\,\mathrm{curv}_j(t_i)}{\partial\theta_k}\frac{\partial\,\mathrm{curv}_j(t_i)}{\partial\theta_l}
    \;+\; \mathrm{curv}_j(t_i)\,\frac{\partial^2\,\mathrm{curv}_j(t_i)}{\partial\theta_k\partial\theta_l}
  \right].
\]
The second term needs \emph{second}-order sensitivities
($\partial^2 Y/\partial\theta_k\partial\theta_l$) -- more expensive, and not
always available. The standard fix (the same one \code{normL2()} itself
uses by default, and the same one behind the Fisher-information
construction discussed previously for the FIM/L2 regularization) is the
\textbf{Gauss--Newton approximation}: drop the second term entirely and
keep only the first, which is a sum of \emph{squares} of first derivatives:
\begin{equation}
  \frac{\partial^2 R}{\partial\theta_k\partial\theta_l} \;\approx\;
  2C\sum_j w_j \sum_i \frac{\partial\,\mathrm{curv}_j(t_i)}{\partial\theta_k}\frac{\partial\,\mathrm{curv}_j(t_i)}{\partial\theta_l}.
  \label{eq:gn}
\end{equation}
\begin{lstlisting}
pen_hessian <- pen_hessian + 2 * C * weight[j] * (t(Jc) %*% Jc)
\end{lstlisting}
\code{t(Jc) \%*\% Jc} computes exactly the $(\theta\times\theta)$ matrix
whose $(k,l)$ entry is $\sum_i \mathrm{Jc}_{i,k}\,\mathrm{Jc}_{i,l}$ --
Equation~\eqref{eq:gn}, for every pair $(k,l)$ at once. Two reasons this
approximation is preferred, not just tolerated: (1) it is
\textbf{guaranteed positive semi-definite} (any matrix of the form
$J^\top J$ is), which keeps the optimizer numerically well-behaved, whereas
the dropped exact term has no such guarantee; (2) it avoids computing
second-order sensitivities at all, which is both cheaper and simpler.

\section{Wiring it into dMod: a wrapper, not a \code{+}-term}

dMod's usual idiom composes penalties with \code{+}
(\code{obj + constraintL2(...)}), and each summand independently calls the
model. That would be wasteful here: the penalty needs the model's
\emph{full predicted trajectory}, so a separately-composed penalty term
would re-solve the ODE a second time on every single objective evaluation.

Instead, \code{curvatureConstraint()} builds a \textbf{wrapper}: it calls
the base objective (\code{normL2}) \emph{once}, and reads the prediction
that call already computed and cached
(\code{attr(base\_out, "env")\$prediction}) before adding the penalty on
top:
\begin{lstlisting}
wrapped <- function(..., fixed = NULL, deriv = TRUE, deriv2 = FALSE, env = NULL) {
  base_out <- base_obj(..., fixed = fixed, deriv = deriv, deriv2 = deriv2, env = env)
  pred <- attr(base_out, "env")$prediction[[condition]]
  ...   # build pen_value / pen_gradient / pen_hessian from `pred`
  combined <- objlist(value = base_out$value + pen_value, ...)
  attr(combined, "env") <- attr(base_out, "env")
  combined
}
class(wrapped) <- class(base_obj)
attr(wrapped, "conditions") <- attr(base_obj, "conditions")
attr(wrapped, "parameters") <- attr(base_obj, "parameters")
\end{lstlisting}
One ODE solve (with sensitivities) serves both the likelihood \emph{and}
the curvature penalty. The last three lines matter more than they look:
without them, \code{wrapped} is just a plain R function, and dMod's own
internals -- specifically \code{mstrust()}, which calls a generic function
\code{mname()} that dispatches on class -- fail with ``no applicable
method for `mname' applied to an object of class `function'\,''. Copying
the base objective's class and attributes is what makes the wrapper behave,
to the rest of dMod, exactly like an ordinary objective function.

\section{Getting states ``for free'': \code{attach.input}}

The penalty is meant to cover \emph{every} state, not just the measured
observable(s) -- but only the observable was ever fit to data. Two
observation functions are built from the same underlying model:
\begin{lstlisting}
g      <- Y(observables, x, ..., attach.input = FALSE)  # observables only
g_full <- Y(observables, x, ..., attach.input = TRUE)   # + every raw state
\end{lstlisting}
Internally, computing $Y(t)=g(X(t))$ always requires solving for the full
state trajectory $X(t)$ first -- the observable is a cheap algebraic
function applied \emph{afterward}. \code{attach.input} only controls
whether that already-computed state trajectory (and its already-exact
sensitivities) is kept in the output or discarded; it costs one array
append, not a second ODE solve. \code{g} (states discarded) is used only to
generate the simulated dataset, so the simulated ``data'' never
accidentally includes a measurement of a state that was never meant to be
observed; \code{g\_full} (states kept) is used to build the fitting
objective, so the curvature penalty has every state available.

\section{The dense-grid trick}

\begin{lstlisting}
times_fit <- seq(0, 80, by = 0.1)
obj <- normL2(data, g_full * x * p, times = times_fit)
\end{lstlisting}
Fitting happens on a dense, \emph{evenly-spaced} grid, even though the
actual data (\code{times\_data <- seq(0, 80, by = 8)}) is far sparser. This
does not change the fitted likelihood at all: \code{normL2()}'s internal
scoring only reads the rows matching real data timestamps; every other row
in the dense grid is simply extra, unused-by-the-likelihood output that the
ODE solver was asked to also produce. The same expensive integration that
computes the likelihood therefore produces, as an unavoidable side effect,
exactly the dense, uniformly-spaced array the curvature penalty needs
(\S4's uniform-grid requirement) -- at no extra integration cost.

\section{The demonstration model, and why it needs regularizing}

\subsection{A feedforward relay -- structurally incapable of oscillating}
The model (\file{derivative\_constraint\_relay.R:105-132}) is a strict
one-way relay: a decaying receptor pulse drives \code{S1}, which drives
\code{S2}, which drives a two-step \code{S3} tier, which drives a two-step
\code{S4} tier, which drives \code{TF}$\to$\code{TFP}$\to$\code{TFPnuc}.
Every reaction only ever regulates something \emph{downstream} of itself;
nothing feeds back. A chain of one-way handoffs like this can only rise and
fall once per stage -- there is no mechanism available to create the
ringing/overshoot that a feedback loop (like the MAPK cascade's inhibitory
loop in \file{derivative\_constraint.R}) can produce. This isolates the
phenomenon of interest from oscillation entirely.

\subsection{Practical non-identifiability, engineered by observation design}
Only the very last quantity, \code{obsTF} $=$
\code{scale\_TF}$\times$\code{TFPnuc}, is ever measured, and only every 8
time units, with noise. Every one of the other 16 states -- in particular
\code{S2}/\code{S2P}, four reaction steps upstream -- is never observed.
Mathematically, the model is still \emph{structurally} identifiable (with
infinite, perfect, fully-observed data, every parameter would be
recoverable); but with \emph{this} sparse, single-readout, noisy dataset,
many different upstream trajectories can produce statistically
indistinguishable downstream signals. That gap between ``identifiable in
principle'' and ``identifiable from this actual data'' is exactly what
\emph{practical} non-identifiability means.

\subsection{A real bug that is worth understanding: signal-to-noise}
The first version of this model used \code{scale\_TF = 1}, giving
\code{obsTF} a true range of about $0$--$0.0009$ -- thirty times
\emph{smaller} than the assumed noise $\sigma=0.03$. With noise that much
bigger than the entire signal, the ``data'' is close to pure noise, and
\emph{any} reasonable fit looks about as good as any other -- for the wrong
reason (no real signal to speak of), not because of the intended
mechanism (sparse sampling of a real signal). This was caught empirically
by printing the true trajectory's actual range and comparing it to
$\sigma$, and fixed by setting \code{scale\_TF = 1000}, bringing
\code{obsTF}'s peak to $\approx 0.92$ -- a sane few-percent relative noise
level. The general lesson: whenever you fix a noise level by analogy to a
different model, check that it is actually small relative to \emph{this}
signal's own dynamic range, or you are not testing what you think you are.

\subsection{What actually happens without regularization}
A 100-start multistart's single best-likelihood (converged) fit reconstructs
\code{S2P} with roughness $\approx 8.6\times$ the true trajectory's own
roughness, and RMSE $\approx 0.15$ against the true trajectory -- while
paying almost nothing in likelihood for it. Several other upstream/interior
quantities are worse still (e.g.\ \code{S3PP}: RMSE $0.95$, roughness
$102\times$ truth). This is the mechanism working exactly as engineered:
the optimizer is free to overfit noise into the \emph{unobserved} parts of
the model, because nothing in the sparse downstream data forbids it.

\section{Checking against ground truth (only possible because this is simulated)}

Because the data is simulated, the \emph{true} trajectory is known --
something you would never have with real data. \code{roughnessAll()} and
\code{rmseVsTruth()} (\file{derivative\_constraint\_relay.R:208-217}) use it
purely as an independent check, never fed into either fit:
\begin{lstlisting}
roughnessAll <- function(pars) {
  pred <- (g_full * x * p)(times_fit, pars)[[1]][, cols_all, drop = FALSE]
  n <- nrow(pred)
  curv <- (pred[3:n,] - 2*pred[2:(n-1),] + pred[1:(n-2),]) / hstep^2   # Eq. (4)
  sqrt(colMeans(curv^2))                                               # RMS curvature
}
rmseVsTruth <- function(pars) {
  pred <- (g_full * x * p)(times_fit, pars)[[1]][, cols_all, drop = FALSE]
  sqrt(colMeans((pred - true_full[, cols_all, drop = FALSE])^2))        # standard RMSE
}
\end{lstlisting}
Applying the curvature penalty (\code{rel\_strength = 0.1}) moves almost
every quantity toward the truth on \emph{both} metrics at once, at a
likelihood cost of $0.004$ (raw likelihood $-54.675 \to -54.671$) while the
total curvature penalty itself drops from $5.42$ to $0.07$ ($\sim$73$\times$)
-- e.g.\ \code{S3PP} RMSE $0.95 \to 0.004$ (243$\times$ better), \code{TFP}
RMSE $0.049\to0.0002$ (300$\times$ better). \code{RecA}/\code{Rec} are the
one exception, and it is a ratio artifact rather than a real failure:
\code{RecA} is a pure single-exponential decay with near-zero true
curvature, so even a tiny absolute roughness inflates the \emph{ratio}
against that near-zero reference.

\section{The strength sweep: reading the tradeoff curve}

\code{rel\_strength = 0.1} was picked by hand. \file{Part 11} of the script
repeats the whole construction across
$\mathrm{rel\_strength}\in\{0.01,0.1,1,10,100\}$ (a log-grid, reusing the
already-computed $0.1$ run rather than refitting it), and pools baseline +
all five strengths into one combined waterfall
(capped to the best 400 of 600 fits by value, using the same
\code{topByValue()} helper introduced for \file{fim\_l2\_constraintL2.R} --
see that companion note -- so one badly-diverged fit can't stretch the
y-axis and hide everything else).

The actual result traces out a textbook bias/cost tradeoff curve:
\begin{center}
\begin{tabular}{@{}lrrrr@{}}
\toprule
\code{rel\_strength} & raw likelihood & curvature penalty & RMSE(\code{S2P}) & roughness(\code{S2P})/truth \\
\midrule
0.01  & $-54.675$ & 0.011  & 0.0306 & 0.437 \\
0.1   & $-54.671$ & 0.074  & 0.0322 & 0.091 \\
1     & $-54.616$ & 0.533  & 0.0326 & 0.018 \\
10    & $-54.262$ & 4.139  & 0.0357 & 0.003 \\
100   & $-45.970$ & 30.661 & 0.0365 & 0.003 \\
\bottomrule
\end{tabular}
\end{center}
Raw likelihood is essentially flat from $0.01$ through $10$ (a change of
$0.4$, against a baseline around $-55$), then collapses sharply at $100$
($-45.97$) -- a clear, numerically visible signature that the penalty has
started to actively fight the data fit rather than just cleaning up
unconstrained roughness. Meanwhile \code{S2P}'s RMSE against truth barely
moves across the whole range (already fixed by even the weakest
regularization), while its roughness keeps falling smoothly. The practical
reading: anywhere from $0.01$ to $10$ is a reasonable choice here; $100$ is
visibly over-regularized, and the likelihood column is what tells you so,
not just visual inspection of a plot.

\section{Guide to the saved plots}

All plots live in \file{scripts/plots/relay\_derivative\_constraint/}.
\begin{itemize}[nosep]
  \item \file{...\_baseline\_waterfall.pdf}, \file{...\_curvature\_waterfall.pdf} --
  one multistart's sorted-by-value fits (\S16's per-run waterfalls), unregularized
  vs.\ \code{rel\_strength=0.1}.
  \item \file{...\_observable.pdf} -- truth / baseline / curvature-regularized,
  for the one thing actually fit to data (\code{obsTF}).
  \item \file{...\_all\_states.pdf} -- the same three trajectories for
  \emph{every} state -- this is where the smoothing effect on the unobserved
  upstream tiers becomes visible directly.
  \item \file{...\_str\{0.01,1,10,100\}\_waterfall.pdf} -- the individual
  per-strength waterfalls from the sweep.
  \item \file{...\_strength\_sweep\_waterfall.pdf} -- all six runs (baseline
  + 5 strengths) pooled into one waterfall, colored by group.
  \item \file{...\_strength\_sweep\_observable.pdf},
  \file{...\_strength\_sweep\_all\_states.pdf} -- truth, baseline, and the
  best fit at every strength, overlaid.
\end{itemize}

\section{Summary: the shape of the whole construction}

\begin{enumerate}[nosep]
  \item Approximate $\ddot Y$ by a central second difference in time
  (Taylor series, \S4) -- the \emph{only} approximation in the whole
  construction.
  \item Average its square per quantity, self-normalize by each quantity's
  own reference curvature at a baseline fit (\S6), scale the total by a
  fraction of the baseline likelihood (\S7).
  \item Get the exact parameter-gradient of that discretized penalty by
  applying the \emph{same} fixed linear operator to the model's own exact
  sensitivities, since a fixed linear time-operator and
  $\partial/\partial\theta$ commute (\S9) -- no finite-differencing of
  $\theta$, ever.
  \item Approximate the Hessian by Gauss--Newton, for the same reasons
  \code{normL2()} itself does (guaranteed positive semi-definite, no
  second-order sensitivities needed) (\S10).
  \item Wrap, rather than \code{+}-compose, so one ODE solve serves both
  the likelihood and the penalty (\S11); reuse the intermediate state
  trajectory that observation functions compute anyway (\S12); fit on a
  dense grid that costs nothing extra because the likelihood only reads
  the sparse subset it actually needs (\S13).
\end{enumerate}

\end{document}
